\section{OHIE \& Causal Forests}

Having solidified our understanding of the mechanics of causal forests, we are now ready to apply them to the Oregon Health Insurance Experiment (OHIE) data. \cite{taubmanMedicaidIncreasesEmergencyDepartment2014} have estimated the average treatment effect of Medicaid on emergency department visits to be 0.408 with a p-value lower than 0.001. This result is, therefore, statistically significant and it means that a person under Medicaid is expected to have 0.408 more visits to the emergency department  than what she would have had with no insurance. We reproduce this result in the analysis code under \url{https://github.com/denismaciel/causal-forest-ohie/tree/master/R}. With causal forest, however, we will output one treatment effect estimate per person and by doing that, try to capture the eventual heterogeneity in the data.

\subsection{Variable Selection}

The datasets of OHIE combined have more than 200 variables. This is a high number of dimensions to analyze and is probably well beyond any researcher's ability to consider every possible interaction between those variables\footnote{If we allow for a causal tree to be constructed with only 2 variables, there $40,000 (200 \times 200)$ different trees one could build} and how these interactions might affect the number of visits a person would have with and without Medicaid. 

Here we will consider only 40 out of the over 200 different variables. An exhaustive list of the variables used to train our causal forest can be found in Table \ref{tab:features} in the Appendix. We have removed all variables that were collected \textit{after} the lottery. The reason for that is that policymaking is one of the relevant areas causal forests can potentially be applied to. To include information that only became available after the treatment has been applied would prevent the results from being used in many policy applications, since the policy decision itself is often about to whom treatment should be applied and (for obvious reasons) one doesn't have access to data collected after the application of the treatment at the time such a decision is made.

We have also excluded variables whose meaning not clear from the data dictionaries provided by the authors of the original study. For example, there are some variables related to previous participation in other welfare programs offered by the state of Orgeon. Although these variables might as well improve our causal models, it adds a bit of overhead to the reader and to the analysis itself. For that, we opted to leave them out.

\subsection{Modeling \& Code}

Before modeling, we have randomly split the data into a train and a test set. The train set contained 80\% of the observations (19,674) while the test set contained the remaining 20\% (4,918). We wanted to have an as close as possible overview of how the model would predict the causal effect for new, completely unseen data. Even though the splitting-estimation-sample procedure, which is already embedded in the causal tree algorithm, ensures that not the same data is used to grow the tree and to estimate the causal effects in the nodes, we go a step further and set aside 20\% of the data, which will remain completely unseen by our causal forest. By doing so, we hope to have a better understanding of how the model will perform on out-of-sample data, which, again, is very important for public policy.

For the modeling itself, we used the \texttt{grf} R package\footnote{The package is open-source and can be found at \url{https://github.com/grf-labs/grf}.}, which is maintained by the authors of the Generalized Random Forest paper. We have used the \texttt{instrumental\textunderscore forest}, which is a generalization of \texttt{causal\textunderscore forest} that accept an instrument together with the independent variables, the treatment and the dependent variable. The causal forest was composed of 2,000 trees. This is the number of trees recommended by the authors if you wish to also estimate confidence intervals along with the treatment effects.

\subsection{Results}

After training the causal forest, we use it to predict the causal effect of Medicaid on the number of visits to the emergency department on the observations in the test set. Figure \ref{fig:rank} shows the estimated causal effects for every observation in the test set ranked by the size.

\begin{figure}
    \includegraphics[width=\linewidth]{ete_p_rank.png}
  \caption{Ranked Treatment Effects}
  \label{fig:rank}
\end{figure}

The estimated individual causal effects range from -4.8 to 4.9 with two extreme outliers at -9.8 and 5.4. Compared to the average treatment effect estimate of the original paper (0.408), it seems to indicate that there might be in fact heterogeneity in the effects and that enrollment in Medicaid could increase emergency department visits for some people and decrease it for other people.

Here we are visualizing almost 5 thousand data points and it's difficult to identify where the majority of the observations concentrate: are there more people with negative or positive treatment effects? Figure \ref{fig:histogram} answer this question by making clear that most of the predictions concentrate on the right side of the line of real numbers. This is in line with the original findings and is indicative that Medicaid leads to an increase in emergency department usage for most people.

\begin{figure}
    \includegraphics[width=\linewidth]{ete_histogram.png}
  \caption{Treatment Effect Histogram}
  \label{fig:histogram}
\end{figure}

As mentioned, one of the greatest features of causal forests is not only to provide estimates on the heterogeneity of treatment effect but also that these estimates are asymptotically normally distributed. Armed with that, we can construct confidence intervals around our point estimates. In Figure \ref{fig:rank_interval}, we have sampled 200 observations from the test set and displayed the predicted treatment effect estimates accompanied by their respective confidence intervals. There is a great amount of uncertainty on the estimates: even for the highest positive estimates, we cannot reject the null hypothesis with 95\% confidence that they are greater than zero.

\begin{figure}
    \includegraphics[width=\linewidth]{ete_p_rank_interval.png}
  \caption{Estimated Treatment Effects \& Confidence Intervals}
  \label{fig:rank_interval}
\end{figure}

\subsection{Feature Importance}

Although many machine learning models are known for being difficult to interpret (black box models),\footnote{There have been many efforts lately to make understandable the inner workings of machine learning models. For a good (and free) overview of the techniques on how to interpret these models, see \cite{molnarInterpretableMachineLearning2019}.} feature importance is an interesting metric that allows us to identify which variables are the most relevant for random forests when making decisions.

With so many variables been fed into the causal forest, it is fairly unlikely that all independent variables are equally relevant for estimating causal effects. The feature importance metric assigns to each independent variable a numeric value that represents how relevant it is for the random forest when making its predictions (\cite{hastieElementsStatisticalLearning2009}). Feature importance is calculated according to the following description:

\begin{itemize}
  \item For each node of every tree in the causal forest, determine the variable that splits it. 
  \item For each independent variable, sum up the resulting improvement in the loss function (here: the increase in the difference between treatment and control related to emergency department usage) whenever this variable is used as the splitting criteria.
  \item Normalize the results in such a way that every feature importance lies between 0 and 1 and their sum equals to 1.
\end{itemize}

Table \ref{tab:feature_importance} shows the feature importance for the ten most relevant variables. The most important variable is ``Number of primary care treatable Emergency Department visits'' with 16.66\%. This type of visits encompasses the conditions that require immediate care but could have been treated in a primary care setting. The second and third most important variables are ``Sum of total charges'' and ``Sum of total Emergency Department charges'', which are the amount expended with hospitals in general and with emergency department only, respectively. Except for ``Age in 2009'', all the eight most important variables are related to previous visits to emergency department. Feature importance, however, does not tell us how each feature affects the number of emergency department visits. 

\begin{table}[!ht]
  \centering
  \begin{tabular}{rll}
    \hline
   & Feature & Importance \\ 
    \hline
    1 & num\_epct\_pre\_ed & 16.66\% \\ 
    2 & charg\_tot\_pre\_ed & 14.87\% \\ 
    3 & ed\_charg\_tot\_pre\_ed & 14.51\% \\ 
    4 & num\_edcnnp\_pre\_ed & 10.89\% \\ 
    5 & num\_ne\_pre\_ed & 10.81\% \\ 
    6 & num\_edcnpa\_pre\_ed & 7.08\% \\ 
    7 & age\_2009 & 4.31\% \\ 
    8 & num\_unclas\_pre\_ed & 2.58\% \\ 
    9 & any\_hosp\_pre\_ed\_Yes & 2.46\% \\ 
    10 & any\_chron\_pre\_ed\_Yes & 2.44\% \\ 
     \hline
  \end{tabular}
  
  \caption{\label{tab:feature_importance} Feature importance of most important variables}
  \end{table}

\newpage
